\usepackage{amscd}

\usepackage[utf8]{inputenc}
\usepackage[T1]{fontenc}

\usepackage{tikz}
\usepackage{pgffor}
\usepackage{tkz-euclide}
\usepackage{tikz-cd}
\usepackage{tikz-3dplot}
\usepackage{pgfplots}

\usepackage{amsmath}
\usepackage{esint}
\usepackage{amsfonts}
\usepackage{mathtools}
\usepackage{amsthm}
\usepackage{amssymb}
\usepackage{mathrsfs}
\usepackage{breqn}
\usepackage{upgreek}
\usepackage{derivative}

\usepackage{epstopdf}
\usepackage{etoolbox}
\usepackage{fancyhdr}

\usepackage{subcaption}
\usepackage[colorlinks=true]{hyperref}
\usepackage{url}
\usepackage{booktabs}
\usepackage{nicefrac}
\usepackage{microtype}
\usepackage{cleveref}
\usepackage{lipsum}
\usepackage{graphicx}
\usepackage{doi}

\usepackage{todonotes}

\captionsetup[figure]{font=small}

\setcounter{tocdepth}{2}

\theoremstyle{plain}
\newtheorem{note}{Note}
\newtheorem{theorem}{Theorem}
\newtheorem{acknowledgement}{Acknowledgement}
\newtheorem{algorithm}{Algorithm}
\newtheorem{axiom}{Axiom}
\newtheorem{case}{Case}
\newtheorem{claim}{Claim}
\newtheorem{conclusion}{Conclusion}
\newtheorem{conjecture}{Conjecture}
\newtheorem{property}{Property}
\newtheorem{corollary}[theorem]{Corollary}
\newtheorem{criterion}{Criterion}
\newtheorem{definition}[theorem]{Definition}
\newtheorem{example}[theorem]{Example}
\newtheorem{exercise}{Exercise}
\newtheorem{lemma}[theorem]{Lemma}
\newtheorem{notation}[theorem]{Notation}
\newtheorem{problem}[theorem]{Problem}
\newtheorem{proposition}[theorem]{Proposition}
\newtheorem{remark}[theorem]{Remark}
\newtheorem{solution}{Solution}
\newtheorem{summary}{Summary}
\numberwithin{equation}{chapter}
\numberwithin{theorem}{chapter}
\newcommand{\thmref}[1]{Theorem~\ref{#1}}
\newcommand{\secref}[1]{\S\ref{#1}}
\newcommand{\lemref}[1]{Lemma~\ref{#1}}

\newcommand\bra[1]{\left\langle #1 \right\rangle}
\newcommand\brak[1]{\left\langle #1 \right\rvert}
\newcommand\braket[1]{\left\langle #1 \right\rangle}

\newcommand\sfrac[2]{#1/#2}
\newcommand\aand{\quad\text{and}\quad}

\newcommand\B{\mathcal{B}}
\newcommand\C{\mathcal{B}}
\newcommand\F{\mathcal{F}}
\newcommand\K{\mathcal{K}}
\renewcommand\L{\mathcal{L}}
\newcommand\M{\mathcal{M}}

\renewcommand\S{\mathbb{S}}
\newcommand\T{\mathbb{T}}
\newcommand\RP{\mathbb{RP}}
\newcommand\N{\mathbb{N}}
\renewcommand\P{\mathbb{P}}
\newcommand\R{\mathbb{R}}
\newcommand\Z{\mathbb{Z}}

\newcommand\VR{\mathrm{VR}}

\usetikzlibrary{
  shadings,
  intersections,
  angles,
  quotes,
  calc,
  positioning,
  3d,
  perspective,
  arrows,
  arrows.meta,
  patterns,
  decorations.markings,
  bending,
  decorations.pathreplacing,
  calligraphy,
  backgrounds,
  shapes.arrows,
}

\tikzoption{canvas is xy plane at z}[]{%
  \def\tikz@plane@origin{\pgfpointxyz{0}{0}{#1}}%
  \def\tikz@plane@x{\pgfpointxyz{1}{0}{#1}}%
  \def\tikz@plane@y{\pgfpointxyz{0}{1}{#1}}%
  \tikz@canvas@is@plane}

\newcommand\graphslopefield{
  \pgfmathsetmacro{\hx}{(\xmax-\xmin)/\nx}
  \pgfmathsetmacro{\hy}{(\ymax-\ymin)/\ny}
  \foreach \i in {0,...,\nx}
  \foreach \j in {0,...,\ny}{
    \pgfmathsetmacro{\yprime}{f({\xmin+\i*\hx},{\ymin+\j*\hy})}
    \draw[black,shift={({\xmin+\i*\hx},{\ymin+\j*\hy})}]
    (0,0)--($(0,0)!2mm!(.1,.1*\yprime)$);
  }

  \draw[->] (\xmin-.5,0)--(\xmax+.5,0) node[below right] {$x$};
  \draw[->] (0,\ymin-.5)--(0,\ymax+.5) node[above left] {$y$};
}

\tikzset{derivative/.style={color=gray,mark=none,line width=0.5pt,solid}}
\tikzset{asymptote/.style={color=gray,mark=none,line width=1pt,<->,dashed}}
\tikzset{soldot/.style={color=black,fill=black,only marks,mark=*}}
\tikzset{holdot/.style={color=black,fill=white,only marks,mark=*}}

\tikzset{>=stealth}
\tikzset{->-/.style={decoration={markings,mark=at position .5 with {\arrow{>}}},postaction={decorate}}}

\usepgfplotslibrary{
  external,
  colormaps,
  groupplots,
}
\pgfplotsset{width=7cm,compat=1.8}
\pgfplotsset{compat=newest}

\let\div\relax
\DeclareMathOperator\diag{diag}
\DeclareMathOperator\div{div}
\DeclareMathOperator\grad{grad}
\DeclareMathOperator\Id{Id}

\hypersetup{
  urlcolor=blue, citecolor=red,
  pdftitle={},
  pdfsubject={q-bio.NC, q-bio.QM},
  pdfauthor={Hashem A.~Damrah},
  pdfkeywords={},
}

% Symbols
\let\oldlimsymbol\lim\renewcommand\lim{\displaystyle\oldlimsymbol}
\let\to\rightarrow
\let\implies\Rightarrow
\let\impliedby\Leftarrow
\let\iff\Leftrightarrow
\let\epsilon\varepsilon
\let\phi\varphi
\let\tau\uptau

% Page settings
\textheight=8.2 true in
\textwidth=5.0 true in
\topmargin 30pt
\setcounter{page}{1}

% The next 5 line will be entered by an editorial staff.
\def\currentvolume{X}
\def\currentissue{X}
\def\currentyear{200X}
\def\currentmonth{XX}
\def\ppages{X--XX}
\def\DOI{10.3934/fods.2019001}

% The next 3 lines are for the arXiv preprint.
\makeatletter
\renewcommand{\xleftrightarrow}[2][]{\ext@arrow 3359\leftrightarrowfill@{#1}{#2}}
\newcommand{\xdashrightarrow}[2][]{\ext@arrow 0359\rightarrowfill@@{#1}{#2}}
\newcommand{\xdashleftarrow}[2][]{\ext@arrow 3095\leftarrowfill@@{#1}{#2}}
\newcommand{\xdashleftrightarrow}[2][]{\ext@arrow 3359\leftrightarrowfill@@{#1}{#2}}
\def\rightarrowfill@@{\arrowfill@@\relax\relbar\rightarrow}
\def\leftarrowfill@@{\arrowfill@@\leftarrow\relbar\relax}
\def\leftrightarrowfill@@{\arrowfill@@\leftarrow\relbar\rightarrow}
\def\arrowfill@@#1#2#3#4{%
  $\m@th\thickmuskip0mu\medmuskip\thickmuskip\thinmuskip\thickmuskip
   \relax#4#1
   \xleaders\hbox{$#4#2$}\hfill
   #3$%
}

% Abstract environment (for amsbook)
\renewenvironment{abstract}{
  \begin{center}
    \bfseries\abstractname
  \end{center}
  \begin{quote}
}{
  \end{quote}
}

% Differentials
\newcommand\dd[1]{\odif{#1}}
\newcommand\pd[1]{\pdif{#1}}
\newcommand\Dd[1]{\mdif{#1}}

% Main macro
\newcommand\dA{\dd{A}}
\newcommand\dF{\dd{F}}
\newcommand\df{\dd{f}}
\newcommand\dV{\dd{V}}
\newcommand\dm{\dd{m}}
\newcommand\ds{\dd{s}}
\newcommand\dS{\dd{S}}
\newcommand\dSS{\dd{\S}}
\newcommand\dw{\dd{w}}
\newcommand\dt{\dd{t}}
\newcommand\dr{\dd{r}}
\newcommand\drr{\dd{\r}}
\newcommand\dT{\dd{T}}
\newcommand\du{\dd{u}}
\newcommand\dv{\dd{v}}
\newcommand\dx{\dd{x}}
\newcommand\dy{\dd{y}}
\newcommand\dz{\dd{z}}

\newcommand\dxy{\dd{x,y}}
\newcommand\dyx{\dd{y,x}}
\newcommand\dzy{\dd{z,y}}
\newcommand\dyz{\dd{y,z}}
\newcommand\dzx{\dd{z,x}}
\newcommand\dxz{\dd{x,z}}

\newcommand\dxyz{\dd{x,y,z}}
\newcommand\dxzy{\dd{x,z,y}}
\newcommand\dyxz{\dd{y,x,z}}
\newcommand\dyzx{\dd{y,z,x}}
\newcommand\dzxy{\dd{z,x,y}}
\newcommand\dzyx{\dd{z,y,x}}

\newcommand\duv{\dd{u,v}}
\newcommand\dvu{\dd{u,y}}

\newcommand*{\bfrac}[2]{\genfrac{}{}{0pt}{}{#1}{#2}}

\title{Geometric Inverse Problems}
\author[Hashem A. Damrah]{}
\date{\today}

\def\subjclass{Primary 53C21; Secondary 35R30, 53C22}
\def\keywords{Riemannian manifolds, inverse problems, geodesics, metric rigidity, X-ray transform, boundary distance data.}
\def\subjclassname{}
\email{hdamrah@uoregon.edu}

\thanks{}

\newcommand\bs{Berry and Sauer~\cite{1606.02353}}
\newcommand\blfootnote[1]{%
  \begingroup
  \renewcommand\thefootnote{}\footnote{#1}%
  \addtocounter{footnote}{-1}%
  \endgroup
}

\newcommand\codv[3][]{\odv[style-inf-num=\mathrm{D},#1]{#2}{#3}}
